\chapter{Produce and Revise a Report}\label{ch:reports}
\section{Define the reader's question and knowledge}

A report is an explanation for a named reader, built around the question that reader
needs to answer. Start by naming that question, the decision or use, the available
evidence, and the report's status. Then record what the reader can be expected to know.
The session's \code{METHODOLOGY_KNOWLEDGE_BASE.md} (or the equivalent in your project)
should connect terminology, mathematics, coursework, literature, and unfamiliar methods
to the planned explanation \cite[Audience and authority; Design the
  argument]{skill:report-production}.

\textbf{Worked example.} Suppose the code rewrite from Chapter~\ref{ch:code}
needs a report answering: ``What evidence supports the claim that this rewrite
preserves the agreed transformation, and what remains untested?'' The input
is its stage contract, source revisions, baseline and candidate artifacts,
comparison records, and performance evidence when required. The report should
explain the acceptance argument; it need not reproduce every execution log.
This example describes a report design, not a claim that the rewrite passed.

\contract\ Assume the reader understands histograms and the analysis context,
but explain any unfamiliar comparison method before using it to accept or reject
agreement. Prefer methods already grounded in the reader's coursework or mapped sources.
A custom diagnostic needs its motivation, definition, limitation, and status
stated explicitly. An unfamiliar method cannot become self-explanatory merely
because it produces a plot \cite[Design the argument]{skill:report-production}.

\section{Inspect the whole artifact and choose its form}

For an existing report, read its complete source, outline, generated evidence, and
rendered PDF before editing. Choose a local patch when the organization still answers the
question; a structural revision when the argument needs consolidation or reordering; and
a rewrite when the organizing model no longer serves the reader. Repetition, appended
findings, and chronology crowding out current conclusions are reasons to reconsider the
structure \cite[Inspect before editing]{skill:report-production}.

Choose the lightest form that gives the reader a clear route. A memo can answer one
bounded comparison; an article can develop one continuous study; a book can carry several
dependent questions and durable reference material. For the illustrative rewrite
comparison, a short article may suffice if there is one argument and limited evidence. A
broader analysis guide may need chapters. Scope and dependency determine the form, not
the template already open \cite[Choose the artifact]{skill:report-production}.

Before drafting, assign each piece of material a role. In our example:
\begin{enumerate}
  \item The opening states the question, bounded coverage, current finding, and review status.
  \item The method explains the frozen contract and how the comparisons test it.
  \item The evidence follows stage order and makes differences inspectable.
  \item Interpretation explains what the evidence supports and what it cannot.
  \item An appendix preserves exhaustive comparisons, commands, and provenance needed to
        reproduce the check.
\end{enumerate}
This is an example outline. The canonical workflow requires a coherent argument,
not these exact headings \cite[Design the argument]{skill:report-production}.

\section{Make every important claim independently checkable}

The bibliography is part of the reader's working apparatus. A reader who can identify a
source, locate the relevant passage, and return to the report's use of it checks the
report far more easily. New substantial reports and substantial rewrites therefore use
shared BibTeX-format \code{.bib} databases, \code{biblatex} with Biber, linked citations,
and bibliography back-references \cite[Bibliography and reader-verifiable
  pinpoints]{skill:report-production}.

Keep the shared library in one directory. External literature has one owner in one file;
maintained project references use a second; dated evidence is grouped by research topic
in further files \cite[Register the source you actually inspected]{skill:library}. Each
source has one stable key and the same descriptive short name across reports. A study
entry may identify a defined evidence collection: the citation then names the particular
member record, while its manifest retains input identities and hashes. A rerun must not
silently replace the version used as evidence by an earlier report.

Three pieces of information do different jobs:
\begin{description}
  \item[Source identity] The bibliography supplies checked author or organization, title,
        edition/version, date, stable identifier where available, and access restriction. A
        repository source instead needs its exact relative path, owning role, and inspected
        revision or date.
  \item[Claim pinpoint] The citation beside the statement names the source's printed page,
        section, equation, table, figure, or record supporting that claim.
  \item[Return route] The bibliography's ``cited on'' pages lead back to this report's uses of
        the source. They are report pages, not source pinpoints.
\end{description}
These distinctions also apply to the report's tables and captions
\cite[Source traceability and corroboration; Bibliography and reader-verifiable pinpoints]{skill:report-production}.

A citation names the source that owns the claim. The hierarchy is:
\begin{itemize}
  \item \textbf{Evidence claims} cite a machine-generated source (runtime output, log,
        executed configuration, code \code{file:line} at a commit) or a human-generated one
        (group documentation, wiki pages, papers, the project's maintained manual). Small cited
        artifacts are copied with their SHA-256 into the session's \code{evidence/} folder,
        because scratch storage is not durable.
  \item \textbf{Session records} (brief, state, result, review and worker records) are cited
        only for the owner's decisions and for review acceptance, never as evidence of a measured
        or implemented fact.
  \item \textbf{Recommendation and deprecation claims} cite the entry in the project's maintained manual or
        register, when one exists. It owns their external citations and records where
        documentation and code disagree. A missing entry is added to the register first.
  \item \textbf{Physics reasoning and tool usage} cite human-facing documentation first: official
        group pages, wiki topics with their revision, tool documentation, and papers. A merge request, commit or \code{file:line} is cited as well
        when that documentation does not point to the change, or when a report explains
        exactly how a tool behaves.
\end{itemize}
Citation and cross-reference links are blue (\code{colorlinks} with \code{ctpBlue}), never
hidden. Reports use descriptive names instead of labels borrowed from other projects
\cite[Source traceability and corroboration; Source hierarchy; Bibliography and
  reader-verifiable pinpoints; Writing character]{skill:report-production}.

\hook{reports-sources}

The following is citation syntax, with placeholder keys and locators to replace after
inspecting the actual sources:
\begin{lstlisting}
\cite[p.~23, sec.~3.9]{source-key}
\cites[pp.~10--11]{first-key}[sec.~2]{second-key}
\end{lstlisting}
Each source in a grouped citation receives its own pinpoint. A bibliography entry's
\code{pages} field usually describes the publication's extent or article number; it does
not identify the passage supporting every claim. For a paper, check the printed
pagination of the cited edition. For Markdown, cite the exact section heading. If a
source is unavailable, mark the missing check rather than guessing a page
\cite[Bibliography and reader-verifiable pinpoints]{skill:report-production}.

In the comparison example, cite the stage record for the agreed transformation, the
precise source revision for implementation, and the comparison artifact for observed
agreement or disagreement. A workflow citation can justify the order of checks; it cannot
substantiate a measured result. A dirty working copy must be identified as such rather
than represented as the published base commit \cite[Source traceability and
  corroboration; Bibliography and reader-verifiable pinpoints]{skill:report-production}.

If a figure includes a source PDF panel, record the caption's printed page and figure
identifier separately from the extraction command's \code{page=} index. Front matter can
make these numbers differ. For derived figures, identify both the input evidence and the
method when they have different owners. Preserve access restrictions when reproducing
material \cite[Choose form by information structure; Bibliography and reader-verifiable
  pinpoints]{skill:report-production}.

\subsection{Keep the editable source readable}

Format authored LaTeX with \code{latexindent} and the project's \code{latexindent.yaml}:
two-space indentation and approximately 90-character prose lines. Preserve paragraphs,
deliberate equation and table breaks, and code listings. Check formatting before handback
using the \code{check_report_sources.py} script of the report-writing skill; its
\code{--write} option applies the formatting. Generated tables should acquire readable
source through their generator, so the next build retains it.

Use \code{siunitx} for numerical results, quantities and units. For example:
\begin{lstlisting}
\SI{25}{\giga\electronvolt}
\num{3084925}
\SI{0.0193}{\percent}
\unit{\pico\barn}
\SIrange{20}{100}{\giga\electronvolt}
\end{lstlisting}
Use \code{\SI} rather than \code{\qty}, which the \code{physics} package uses for
delimiter sizing. Keep the original precision and uncertainty notation. Dataset
identifiers, hashes, software versions, dates and algebraic constants retain their own
notation. Numeric table columns can use \code{S} alignment when it helps comparison. The
quantity scan flags likely omissions; checking the prose, captions and tables still
requires distinguishing values from identifiers \cite[Readable source and numerical
  presentation]{skill:report-production}.

\section{Build the explanation around useful evidence}

Introduce a table or figure before asking the reader to interpret it. State what it
compares, what is held fixed, and which feature answers the question. For the rewrite
example, a compact stage-comparison table can expose the first mismatch; a labelled
distribution comparison can show its shape. Exhaustive per-file logs belong in supporting
evidence unless a particular entry explains the conclusion \cite[Design the argument;
  Choose form by information structure]{skill:report-production}.

For a personal working report the target is a short time to read and digest, not a short
page count; note style is fine while it stays clear. A recommended order of preference
for explaining something, which you may reorder for your readers, is:
\begin{enumerate}
  \item plots, sketches and diagrams; then commented equations with each term explained;
  \item lists and tables;
  \item prose.
\end{enumerate}
Build the aids first, then write the prose that presents the argument and connects them.
A code or pseudo-code snippet may explain what code does when it is clearer than prose. If
the reader would need to re-read a paragraph to follow it, it should become an aid. A diagram
explains a relationship; it does not by itself establish numerical agreement. Keep source
data separate from presentation when generating repeated tables so prose and LaTeX do not
acquire independently maintained copies of the same facts \cite[Choose form by information
  structure]{skill:report-production}.

Three devices from the report template make the structure visible \cite[Headlines,
  choices and colour]{skill:report-production}:
\begin{description}
  \item[Headline box] The section's takeaway in one or two sentences, first thing in the section,
        so the payoff does not wait until its end.
  \item[Choice box] Every choice for the reader, itemized, with a recommendation (blue), an
        alternative (mauve) and links to the supporting plots or tables.
  \item[Colour] Shaded table rows, marked outliers and coloured list markers where they help the
        eye, each colour with one stable meaning and a textual cue.
\end{description}
When the story jumps, reorganize: a definition or equation sits next to its first use.
Once a question is settled, a new dated revision may retire the presentation that served
it; earlier revisions keep the history \cite[Inspect before editing]{skill:report-production}.

Useful visual conventions are portrait pages, readable spacing, explicit reading routes,
restrained semantic color, and legible code blocks. Give every color a stable labelled
meaning. A green ``evidence'' heading can identify an observed record, but a proposed
expectation must not look like an already verified result. Inspect grayscale rendering
when color carries meaning \cite[Colour convention]{skill:report-production}.

\section{Iterate figures and the plotbook before finishing the companion}

Start from the equation. Each plotted quantity is an equation, numbered in the report or
written in the plot's how-to-read note, and the plot shows it as a function or its terms
as histograms, with the equation's symbols in the labels. The standard grammar of the
field comes first: a distribution with extra panels for a ratio, a pull or a diagnostic
such as the stacked terms of a bin decomposition. A new plot type makes the reader learn
a new grammar, so it must earn its place, and its note says why. Every axis is labelled:
a shared $x$ axis once, under the lowest panel, and each panel its own $y$ axis; a
mathematical expression is preferred to words unless a name reads better. If you maintain
your own plotting library, add a missing reusable capability there, with a test, rather
than drawing around it in a recipe \cite[Design and iterate at two scales; Reuse the
  implementation]{skill:plot-production}.

Begin with the reader's task and the evidence needed to perform it. Coverage and page
count are separate: a dense plotbook makes useful comparisons easy without discarding
unique measurements. Recover the applicable binning, statistical conventions and plotting
machinery, then plan candidate groupings. The first plan cannot establish how a
complicated overlay will read before it is rendered \cite[Design and iterate at two
  scales; Reuse the implementation]{skill:plot-production}.

There are two connected revision loops:
\begin{description}
  \item[Figure.] State the question, choose an encoding, render actual data and test whether a
        series can be traced from its legend through the main and ratio panels. Colors, line
        styles and markers are tools. Keep them when they help decoding; revise them when they
        disappear or compete at final size.
  \item[Plotbook.] Check coverage, group comparisons, compose pages and inspect the whole
        sequence. Decide what rows and columns mean. Bring alternatives together when
        simultaneous viewing helps; separate them when they obscure one another. A sparse page
        with tiny plots calls for reconsidering the composition.
\end{description}
A page-level finding can reopen figure design. Neither a fixed grid nor a
historical source/flavour split is a universal solution. Agents normally iterate
within the confirmed contract. An unresolved scientific choice or a conflicting preference
is settled by a recorded default that the owner reviews later, unless no default is
defensible \cite[Design and iterate at two scales]{skill:plot-production}.

A recommended habit is to read the plotbook before the report, for a first impression, so
that the report's interpretation does not anchor the reading (the Einstellung effect).
The plotbook must therefore be readable on its own \cite[What a plotbook holds;
  Established appearance]{skill:plot-production}:
\begin{itemize}
  \item legends and colour keys sit inside the plot, and axis labels name the compared stages,
        for example $(N_\mathrm{BM}-N_\mathrm{BH})/N_\mathrm{BH}$, instead of a title or an
        explanation block atop the page;
  \item a plot whose reading is not obvious carries a short how-to-read note pointing to the
        report section or numbered equation that explains it, and no task-specific term is left
        undefined;
  \item the plotbook holds plots: a long table such as a per-event longtable does not belong, a
        small selected table may;
  \item a page with a large empty area is rejected: the render gate fails a blank band taller
        than \SI{25}{\percent} of the page height unless an exception is recorded;
  \item once results transfer between validation cohorts, the development cohort and one
        transfer-check summary replace side-by-side cohorts.
\end{itemize}

Prove the densest or least certain group in its actual page placement. Audit scale
choices across every affected family, including tails, zero/signed bins and the
uncertainty envelopes actually drawn. Reserve annotation space after overlays and final
scales: a later axis-limit change can undo earlier clearance. Numerical preservation and
geometry checks support this work but cannot judge whether the comparison communicates
its effect \cite[Scales and final geometry; Render and accept]{skill:plot-production}.

For a companion report, establish a provisional argument early but settle the plotbook's
comparisons and page map before final prose. Explain why each comparison answers the
question; assume domain knowledge, not participation in the agent sessions. Introduce a
population before its ratio and use a short worked example where it removes a difficult
conceptual step. Give recurring conventions one home. If writing exposes missing
evidence, return to the plotbook and recheck affected routes rather than writing around
the gap \cite[Write the companion after settling the plotbook]{skill:report-production}.

After the complete draft exists, read it as one argument. Remove repeated context,
reconcile old and new conclusions, move chronology out of the main explanation, and
define terms before use. If the comparison has changed the conclusion, revise the opening
and interpretation together rather than appending a corrective paragraph at the end
\cite[Design the argument; Writing character]{skill:report-production}.

\section{Use separate review tracks and bounded corrections}

A report can be grammatical and wrong, numerically sound and unsupported by its
citations, or correct in source but unreadable in the PDF. Table~\ref{tab:reviews}
separates these checks. Internal review prepares the artifact for the owner's own
inspection; it does not replace their judgment.

\begingroup
\small
\setlength{\LTcapwidth}{\textwidth}
\begin{longtable}{L{0.17\textwidth}L{0.34\textwidth}L{\dimexpr0.49\textwidth-6\tabcolsep}}
  \caption{Report review tracks and the evidence each supplies
    \cite[Source traceability and corroboration; Tiered review]{skill:report-production}.}
  \label{tab:reviews}\\
  \toprule Track & Question & Required evidence \\ \midrule
  \endfirsthead
  \toprule Track & Question & Required evidence \\ \midrule
  \endhead
  Source & Does the cited passage or record support the actual claim? & Checked
  source identity, edition, pinpoint, and claim; unresolved conflicts recorded. \\ \addlinespace
  Copy and structure & Can the reader follow the complete argument in the
  intended voice? & Independent complete-source review for a substantial
  report; findings resolved by the writer. \\ \addlinespace
  Domain / numerical & Are the definitions, calculations, and interpretation
  justified? & Capable-model review of the complete artifact when conclusions
  depend on such claims, plus required human domain review. \\ \addlinespace
  Build and bibliography & Does the source compile with all references and
  citation data resolved? & LaTeX and Biber logs, no unresolved keys or references,
  and the final generated PDF. \\ \addlinespace
  Rendered pages & Are text, tables, figures, labels, and navigation usable? &
  Complete applicable mechanical checks, then bounded independent inspection of
  new layouts, representative families and flagged pages; reuse accepted evidence. \\
  \bottomrule
\end{longtable}
\endgroup

During drafting, review changed claims and their immediate dependencies. Once bounded
findings are resolved, a substantial report needs the complete source corroboration and
copy gates; apply the complete capable-model gate when physics or numerical conclusions
require it. After a local correction, repeat only the affected checks unless it changes
structure, argument order, source hierarchy, global terminology, document-wide
conventions, or claims outside that scope \cite[Source traceability and corroboration;
  Tiered review]{skill:report-production}.

For example, a wrong source page requires rechecking the passage and affected claim, then
rebuilding its citation. A changed normalization can alter the method, tables, plots,
interpretation, and conclusion; that is a wider domain and source correction. A long path
overflowing a cell is a rendering defect, not automatically a scientific defect. Resolve
each finding at the layer that owns it \cite[Production workflow; Tiered
  review]{skill:report-production}.

Complete coverage can combine prior accepted evidence and bounded rechecks; it need not
repeat the same work in every reviewer session. Programmatic checks cover geometry, blank
bands, printed values, inventory and navigation where applicable. Missing checks are not
assessed, and an unblocked lint result alone is not readiness. Before visual dispatch,
validate the explicit model and page selection with the render gate's preflight. Use an
economical image-capable model for remaining reading tasks. Repagination reopens
navigation and composition; it does not automatically invalidate accepted page-body
readability. Expensive or exhaustive review requires a concrete unresolved problem
\cite[Tiered review]{skill:report-production}.

Give a copy reviewer the complete editable source. For visual delegation, rasterize first
and provide relevant labelled PNG or JPEG contact sheets and selected pages in a scratch
directory; never hand the visual reviewer a PDF. A vision report is evidence about
appearance and needs comparison with the source and intended rendering before it becomes
a claim about correctness \cite[Tiered review]{skill:report-production}.

Give the independent reviewer the reader brief and artifact before the author's verdict.
Ask it to reconstruct the question, findings and supporting comparisons without the
session narrative. It may find a structural problem even when the sentences are correct.
On rendered pages, ask it to trace a series, read the labels and identify the intended
effect. Contact sheets establish overview coverage; individual pages at reading size are
needed for decoding and small text. An unreadable review image means that feature is not
assessed \cite[Review packet; Figure and page tasks; Companion comprehension
  tasks]{skill:visual-review}.

Record the exact artifact, page, observation, reader impact and recheck result. Verify
alleged defects against the image: agreement with criticism does not prove inspection.
Iterate until reading tasks work and material defects are resolved; neither a fixed
number of passes nor repeated cosmetic changes establish quality \cite[Findings and
  revision]{skill:visual-review}.

\section{Build, inspect, and deliver the exact PDF}

The recommended bibliography setup uses \code{backend=biber}, \code{style=alphabetic},
\code{sorting=nyt}, and \code{backref=true}. Use readable, unambiguous short source
labels where they help. A manual build sequence is shown below; \code{latexmk} runs the
same steps and repeats them until the references settle:
\begin{lstlisting}
pdflatex -interaction=nonstopmode -halt-on-error report.tex
biber report
pdflatex -interaction=nonstopmode -halt-on-error report.tex
pdflatex -interaction=nonstopmode -halt-on-error report.tex
\end{lstlisting}
Inspect both logs for missing or duplicate keys, malformed data, unresolved references,
rerun requests, and material layout warnings. Check the rendered citation-to-entry links
and bibliography return links as well as page appearance \cite[Bibliography and
  reader-verifiable pinpoints]{skill:report-production}.

The authored \code{.tex} and \code{.bib} are the reconstructible source. The PDF and
ordinary bibliography/build intermediates remain generated. A clean build proves that the
source can be rendered; the source, copy, domain, and visual reviews establish the
additional evidence needed for delivery \cite[Acceptance checklist; Bibliography and
  reader-verifiable pinpoints]{skill:report-production}.

Hand back the exact PDF and source paths, review evidence, and remaining limits. For a
requested plotbook companion, deliver the matched pair with checked routes; raw logs
remain supporting material. Keep reviewable copies in one collection, delivered as soon
as they are ready: a same-day revision replaces its file, a later date goes alongside,
with every revision's hashes and human-review state in a manifest \cite[Review and
  release]{skill:report-production}. \hook{reports-delivery} Record authorized external
delivery, notification, and source publication separately \cite[Review and
  release]{skill:report-production}.

